%%%
%
% $Autor: Wings $
% $Datum: 2021-05-14 $
% $Pfad: GitLab/CornerBLending $
% $Dateiname: Hints
% $Version: 4620 $
%
% !TeX spellcheck = de_DE/GB
%
%%%



% -------------------------------------------------------
% File: Hints.tex
% Description: General Assembly Hints with Diagram and Table
% -------------------------------------------------------

\chapter{General Assembly Hints}

Before beginning the assembly of the License Plate Detection System, please read the following important guidelines.  
These instructions ensure a safe, efficient, and error-free setup process for the Jetson Nano hardware and its connected components.

% -------------------------------------------------------
% Section: Preparation and Safety
% -------------------------------------------------------
\section*{Preparation and Safety Measures}

\begin{itemize}
	\item \textbf{Work Environment:}  
	Assemble all components on a clean, dry, and non-conductive surface. Ensure proper lighting and ventilation to avoid overheating during operation.
	
	\item \textbf{Electrostatic Discharge (ESD):}  
	Handle the Jetson Nano and camera only by their edges. Use an antistatic wrist strap or touch a grounded metal object before handling any electronics.
	
	\item \textbf{Cable Management:}  
	Arrange all cables neatly and ensure that they are not under tension. Avoid crossing power and data cables where possible.
	
	\item \textbf{Component Orientation:}  
	Check the correct direction of all connectors (micro-SD, USB, HDMI, and power) before insertion. Do not apply excessive force.
	
	\item \textbf{Power Supply:}  
	Use only the official 5 V -  4 A adapter. Always disconnect power before adding or removing components.
\end{itemize}

% -------------------------------------------------------
% Table: Preparation Checklist
% -------------------------------------------------------
\begin{table}[H]
	\centering
	\caption{Assembly Preparation Checklist}
	\begin{tabular}{|L{4cm}|L{9cm}|}
		\hline
		\textbf{Item} & \textbf{Description / Action Required} \\ \hline
		Work Surface & Clean, dry, non-conductive, and well-lit area \\ \hline
		Power Adapter & Verified 5V - 4A supply (official Jetson Nano adapter) \\ \hline
		SD Card & 32 GB card pre-flashed with JetPack 5.1.2 (Ubuntu 20.04) \\ \hline
		Display and Peripherals & HDMI monitor, USB keyboard, and mouse ready \\ \hline
		Camera Module & Microsoft LifeCam Show (USB) connected but protected with lens cover \\ \hline
		Network Connection & Ethernet or Wi-Fi adapter available for first boot \\ \hline
		Static Safety & ESD wrist strap or grounding point nearby \\ \hline
	\end{tabular}
\end{table}

% -------------------------------------------------------
% Section: Ventilation and Safety
% -------------------------------------------------------
\section*{Ventilation and Safety}

\begin{itemize}
	\item Place the Jetson Nano on a hard, flat surface with sufficient airflow around the heat sink.  
	\item Do not operate the system inside a closed enclosure without ventilation holes.  
	\item Avoid touching the heat sink during or immediately after operation, as it may become hot.  
	\item If long-duration tests are performed, use an external cooling fan for optimal performance.
\end{itemize}

% -------------------------------------------------------
% Section: Post-Assembly Verification
% -------------------------------------------------------
\section*{Post-Assembly Verification}

After completing all connections, verify the following points:

\begin{enumerate}
	\item The power LED on the Jetson Nano turns on when the adapter is connected.  
	\item The HDMI display output is visible on the monitor.  
	\item The USB camera is detected and functioning properly.  
	\item All cables and connectors are securely attached.  
	\item No abnormal heat or noise occurs during operation.
\end{enumerate}

% -------------------------------------------------------
% TikZ Diagram: Basic Assembly Layout
% -------------------------------------------------------
\begin{figure}[H]
	\centering
	\begin{tikzpicture}[
		every node/.style={font=\footnotesize, align=center},
		component/.style={draw, rounded corners, fill=blue!10, minimum width=2.8cm, minimum height=1cm},
		connection/.style={->, thick}
		]
		% Components
		\node[component] (nano) {Jetson Nano Board};
		\node[component, above left=1.5cm and 2cm of nano] (camera) {USB Camera};
		\node[component, above right=1.5cm and 2cm of nano] (monitor) {HDMI Monitor};
		\node[component, below left=1.5cm and 2cm of nano] (keyboard) {USB Keyboard};
		\node[component, below right=1.5cm and 2cm of nano] (mouse) {USB Mouse};
		\node[component, below=2cm of nano] (power) {5V Power Adapter};
		
		% Connections
		\draw[connection] (camera) -- (nano);
		\draw[connection] (monitor) -- (nano);
		\draw[connection] (keyboard) -- (nano);
		\draw[connection] (mouse) -- (nano);
		\draw[connection] (power) -- (nano);
		
		% Labels
		\node[below=0.4cm of power] {\textbf{Figure:} Assembly layout for Jetson Nano hardware setup.};
	\end{tikzpicture}
\end{figure}

\vspace{0.4cm}
\noindent
Following these general assembly hints ensures reliable installation, minimizes hardware risks, and provides a stable foundation for further software configuration and testing.

